% !PRINTED 213
\weberpage{213}
%<<CONT
sind, und diese Gleichung kann auch nicht zerfallen, wenn der unabhängigen Veränderlichen $\varkappa^2$ durch Adjunktion der Rationalitätsbereich der rationalen Zahlen erweitert wird. % [sic] word order as printed; expected sense \textquotedbl{}wenn der Rationalitätsbereich ... durch Adjunktion der unabhängigen Veränderlichen ... erweitert wird\textquotedbl{}

Der Grad dieser Gleichung ist $\varphi(n)$, d.\,h. gleich der Anzahl der Modulo $n$ inkongruenten Zahlen $m$, die relativ prim zu $n$ sind.

Es sei nun

\weq{13}{F'(r) = 0} % [sic] The source prints F' here; the subsequent argument consistently uses F.

eine \gesperrt{Galoissche} Resolvente vom Grade $\mu$ der Teilungsgleichung im ursprünglichen Rationalitätsbereich (so daß alle Wurzeln $x_{\nu,\nu'}$ rational durch $r$ darstellbar sind, Bd.\,I, § 152). Diese Gleichung muß nach Adjunktion einer Wurzel von (12) \gesperrt{zerfallen}; denn wäre dies nicht, so würde jede rationale Relation

\weq{14}{\psi(r,\varrho) = 0}

für \gesperrt{alle} Wurzeln von (13) befriedigt sein, und $\psi(t,\varrho)$ wäre durch $F(t)$ (für ein variables $t$) teilbar. Es würde also (14) noch bestehen bleiben, wenn $\varrho$ durch eine andere Wurzel $\varrho^m$ von (12) ersetzt wird. Dies ist aber nach § 62, (12) nicht zutreffend, wenn man an Stelle von (14) eine der \gesperrt{Abelschen} Relationen setzt.

Es sei nun \gesperrt{nach Adjunktion von} $\varrho$

\weq{15}{F(r,\varrho) = 0}

die \gesperrt{Galoissche} Resolvente vom Grade $\nu$ der Teilungsgleichung, so daß $\nu$ nach 2. gleich dem Grade der Gruppe $\mathfrak{B}$ ist. Es ist dann $F(t)$ durch $F(t,\varrho)$ algebraisch teilbar und also wegen der Irreduzibilität von (12) auch durch jede der Funktionen $F(t,\varrho^m)$. Da die Funktionen $F(t,\varrho^m)$ irreduzibel sind, so können nur dann zwei von ihnen einen gemeinsamen Teiler haben, wenn sie ganz identisch sind. Wenn aber

\[
F(t,\varrho) = F(t,\varrho^m)
\]

wäre, dann würde aus jeder Gleichung der Form (14) folgen, daß $\psi(t,\varrho)$ durch $F(t,\varrho) = F(t,\varrho^m)$ teilbar wäre, und es würde folgen, daß in (14) die \gesperrt{Vertauschung} $(\varrho,\varrho^m)$ gestattet ist, was wieder bei den \gesperrt{Abelschen} Relationen nicht zutrifft. Mithin sind die $\varphi(n)$ Funktionen $F(t,\varrho^m)$ alle voneinander verschieden, und $F(t)$ ist durch ihr Produkt teilbar. Der Grad $\mu$ von $F(t)$ ist also wenigstens $= \nu\varphi(n)$. Er kann aber auch nicht höher als $\nu\varphi(n)$ sein, da $\nu\varphi(n)$ der Grad der Gruppe $\mathfrak{A}$ ist, und die
%CONT>>
